%% SciFi-VIS 2026 submission, position paper
%% Uses the official IEEE VGTC conference class (vgtc.cls).
%% See README_OVERLEAF.md for setup steps.

\documentclass{vgtc}                % conference style, final (not double-blind)
%\documentclass[review]{vgtc}       % use if the workshop switches to blind review

\usepackage{times}
\usepackage{mathptmx}
\usepackage{graphicx}
\usepackage{booktabs}
\renewcommand*\ttdefault{txtt}

\onlineid{0}
\vgtccategory{Research}
\vgtcinsertpkg

%% ---------------------------------------------------------------------------
\title{The Lead-Lag Layer Stack. Encoding Temporal\\ Precedence as Spatial Depth}

\author{Jacie Jermier\thanks{jjermier@uwaterloo.ca}\\ %
        \scriptsize University of Waterloo}

\abstract{%
Wastewater surveillance can detect SARS-CoV-2 outbreaks days to weeks before
clinical case counts rise, yet existing public-facing dashboards communicate this
lead time through static line charts and choropleth maps. We describe two
encodings for lead-lag relationships in spatiotemporal data. A \emph{lead-lag
layer stack} renders two co-registered map layers that separate in depth, so
temporal precedence is read spatially. A \emph{paired lead-lag ridgeline} draws
the measured lag between two curves as a span rather than leaving it to the eye.
We demonstrate both in \emph{The Upside Down}, a narrative visualization of
COVID-19 wastewater surveillance framed by \emph{Stranger Things}, and we report
the places where the encodings fail.%
}

\keywords{Lead-lag encoding, spatiotemporal visualization, figurative frames,
science fiction, wastewater surveillance.}

\teaser{
  \centering
  \includegraphics[width=\linewidth]{Figures/figure_1.png}
  \caption{The lead-lag layer stack. Two co-registered map layers tilt to 65
  degrees and separate in Z-space. The lower red layer carries the leading signal,
  wastewater detection, and the upper blue layer the lagging one, clinical cases.
  Flat they coincide. Tilted, the depth between them is the relationship.}
  \label{fig:teaser}
}
%% ---------------------------------------------------------------------------

\begin{document}
\maketitle

\section{Contribution}

Infected individuals shed SARS-CoV-2 RNA into wastewater before symptoms appear
and before they seek testing~\cite{daoust2021,peccia2020}, so outbreaks can be
detected days to weeks earlier than hospital reporting
allows~\cite{kumar2022}. Existing visualizations of this data remain limited to
time-series overlays and choropleth maps. They are analytically useful for
experts but fail to show what makes the data compelling. The virus was detectable
underground long before it showed up in hospitals.

The problem is not specific to COVID-19. Wherever one spatial signal reliably
precedes another across the same geography, a visualization must convey both
where the signals are and which arrived first. We contribute two encodings.

A \textbf{lead-lag layer stack} draws the two signals as transparent,
co-registered planes. Flat they coincide, and the display reads as an ordinary
map. Tilted they separate in depth, and the gap between them is the relationship
(Fig.~\ref{fig:teaser}). Precedence becomes something a viewer reads spatially
rather than infers from two adjacent charts.

A \textbf{paired lead-lag ridgeline} stacks one row per region, each carrying a
filled curve for the leading signal and a dashed curve for the lagging one. The
measured lag is drawn as a labelled span across the row (Fig.~\ref{fig:ridge}),
so the quantity the plot exists to communicate is stated rather than estimated.

\section{Related Work}

Tilt has been used before to move between map representations. Yang et
al.~\cite{tiltmap} transition between choropleth map, prism map and bar chart in
immersive environments, treating tilt as the interaction that swaps one
representation for another. Our tilt does something different. The representation
never changes. Tilting exposes a depth separation that was already carrying the
lead time, so the gesture reveals structure rather than substituting a chart.

The structure itself is familiar in the earth sciences, where subsurface
precursors such as seismic swarms, ground deformation and gas emissions precede
eruptions by a measurable interval~\cite{usgs}. Those systems generally infer the
subsurface from the surface. Wastewater surveillance runs the other way, using a
buried signal to anticipate a surface one, which is why the two-layer arrangement
fits it so directly.

\section{What the Frame Encodes}

The \emph{Stranger Things} reference operates at two levels, and separating them
matters.

The first is structural. The show's premise is a parallel dimension occupying the
same physical space as the visible world, where threats become detectable before
they surface. That premise is not decoration here. It is the encoding, because it
is the same shape as the data relationship. Byrne et al.~\cite{byrne2019} argue
that figurative elements extend a visualization's expressive range beyond
decoration. We would add that a figurative \emph{structure} can carry the encoding
itself, which is a stronger claim than an element carrying meaning alongside one.

The second level is the figurative elements proper, and they do not all do the
same work. Some encode magnitude. Sites are dots whose radius and brightness
follow viral percentile and population served (Fig.~\ref{fig:detection}), and
outbreaks are cracks whose severity combines a state's case intensity with the
wastewater signal at each site, so the places that detected most crack hardest.
Others encode nothing and exist for atmosphere, including the dust in the void
between the layers and the typography.

One element encodes uncertainty, and it is the most useful of the three.
Wastewater coverage is severely uneven. Idaho contributes one sampled plant
serving under three percent of the state, Illinois seventy-one covering
sixty-two. Sites in thinly covered states are therefore drawn dim and unstable,
dropping in and out as the timeline advances, the way the lights do in the show
when something is only just getting through. Coverage becomes legible before any
number is read, and a barely monitored place can no longer pass for a well
monitored one.

\section{Testing the Frame}

Replacing our original manual lead-time estimates with cross-correlation of weekly
wastewater percentile against weekly cases changed the paper twice.

First, the Delta wave had to be removed. CDC NWSS percentile reporting does not
begin until 15 November 2021, so the entire Delta window contains no viral
intensity data. Delta had been our flagship wave and carried most of our figures,
and the visualization had been rendering it confidently the whole time. Nothing in
the design forced the metaphor to break, which is the exact failure we warn about.

Second, our strongest numbers came from our weakest data. Ranked on correlation
alone, Idaho led Omicron by three weeks at $r = 0.99$, on one plant. We now report
a state only when its sampled plants cover at least 15\% of the population across
five or more sites, and every large lead time disappeared under that rule. What
survives is modest (Table~\ref{tab:leads}). Where a lead is measurable it is about
one week, and during BA.5 most adequately covered states fell below that.

This is where the encoding meets its own floor. Our wastewater readings are daily,
but the case counts are weekly, and a lag can never be finer than the coarser of
the two series. Published leads run from roughly two to eight
days~\cite{daoust2021,peccia2020}, which sits at or under seven. So a row reading
under one week has not been shown to lack a lead. It has been shown to have one
this data cannot resolve, and the frame cannot draw a gap the data cannot see.

\begin{table}[t]
  \centering
  \caption{Lead times from cross-correlation of weekly wastewater viral
  percentile against weekly clinical cases, for states where sampled plants cover
  at least 15\% of the population. Weekly case data puts a seven day floor under
  what can be resolved, so $<$1 week means the lead was too short to measure, not
  that none exists.}
  \label{tab:leads}
  \begin{tabular}{llccc}
    \toprule
    State & Wave & Coverage & Lead & $r$ \\
    \midrule
    Illinois      & Omicron & 62\% & 1 week           & 0.90 \\
    Oregon        & Omicron & 51\% & 1 week           & 0.82 \\
    Wisconsin     & Omicron & 23\% & 1 week           & 0.89 \\
    Georgia       & BA.5    & 22\% & 1 week           & 0.92 \\
    Illinois      & BA.5    & 62\% & \textbf{$<$1 wk} & 0.81 \\
    West Virginia & Omicron & 22\% & \textbf{$<$1 wk} & 0.94 \\
    West Virginia & BA.5    & 22\% & \textbf{$<$1 wk} & 0.92 \\
    \bottomrule
  \end{tabular}
\end{table}

\begin{figure}[t]
  \centering
  \includegraphics[width=\columnwidth]{Figures/figure_2.png}
  \caption{Detection phase. Wastewater sites appear as glowing dots before any
  surface cracks. Sites in thinly sampled states render unstable rather than
  confident.}
  \label{fig:detection}
\end{figure}

\begin{figure}[t]
  \centering
  \includegraphics[width=\columnwidth]{Figures/figure_3.png}
  \caption{The paired lead-lag ridgeline. Each row carries a filled wastewater
  curve and a dashed clinical curve. One row per wave is annotated as a worked
  example, chosen because it is the region where the fitted lag and the visible
  peak-to-peak gap agree.}
  \label{fig:ridge}
\end{figure}

We show Appalachia and West Virginia rather than leaving them out
(Fig.~\ref{fig:counter}). In those rows the two curves sit on top of each other
and the layers have nothing to separate. This unevenness is not a flaw in the
visualization but a reflection of the real variability documented by Kumar et
al.~\cite{kumar2022}. The encoding fails in public, which is the strongest
evidence we have that it was carrying something.

\begin{figure}[t]
  \centering
  \includegraphics[width=\columnwidth]{Figures/figure_4.png}
  \caption{The counterexample. With Appalachian highlighted the two curves are
  coincident. No gap, no lead time, nothing for the layers to separate.}
  \label{fig:counter}
\end{figure}

This is the property we would ask of any speculative frame. Its credibility rests
on the cases where it does not hold, and an interface that cannot fail against
real data is closer to concept art, which is valuable but a different
contribution.

\section{Limitations}

The seven day resolution floor is the sharpest limit. Daily case counts would let
the encoding address the range the literature actually reports, and we did not
have them. Sewer transit time and sampling frequency are not
modeled~\cite{kumar2022}, coverage spans Omicron and BA.5 only, the stack needs
WebGL and hardware-accelerated transforms, and we report correlations rather than
causal estimates.

%% ---------------------------------------------------------------------------
\begin{thebibliography}{9}

\bibitem{byrne2019}
L.~Byrne, D.~Angus, and J.~Wiles.
\newblock Figurative frames. A critical vocabulary for images in information
  visualization.
\newblock {\em Information Visualization}, 18(1):45--67, 2019.

\bibitem{daoust2021}
P.~M. D'Aoust et~al.
\newblock Catching a resurgence. Increase in {SARS-CoV-2} viral {RNA} identified
  in wastewater 48\,h before {COVID-19} clinical tests.
\newblock {\em Science of the Total Environment}, 770:145319, 2021.

\bibitem{kumar2022}
M.~Kumar, A.~K. Srivastava, and N.~Srivastava.
\newblock Lead time of early warning by wastewater surveillance for
  {COVID-19}.
\newblock {\em Chemical Engineering Journal}, 441:135936, 2022.

\bibitem{peccia2020}
J.~Peccia et~al.
\newblock Measurement of {SARS-CoV-2} {RNA} in wastewater tracks community
  infection dynamics.
\newblock {\em Nature Biotechnology}, 38(10):1164--1167, 2020.

\bibitem{tiltmap}
Y.~Yang, T.~Dwyer, K.~Marriott, B.~Jenny, and S.~Goodwin.
\newblock Tilt map. Interactive transitions between choropleth map, prism map and
  bar chart in immersive environments.
\newblock {\em IEEE Transactions on Visualization and Computer Graphics},
  26(12):3577--3589, 2020. doi\string: 10.1109/TVCG.2020.3004137

\bibitem{usgs}
U.S. Geological Survey.
\newblock Movement on the surface provides information about the subsurface.
\newblock Volcano Hazards Program. Accessed 1 August 2026.

\end{thebibliography}

\end{document}
